\chapter{Splenectomy}
\section*{What Is This Test or Procedure?}
Splenectomy is surgical removal of the spleen, performed through an open or laparoscopic approach.
\section*{Why This Test or Procedure Is Ordered}
Indications include traumatic rupture, hypersplenism, selected blood disorders, splenic tumors, cysts, and severe infection or infarction.
\section*{How This Test or Procedure Is Performed}
The splenic vessels are controlled, the spleen is freed from surrounding attachments, and it is removed. Laparoscopic surgery is common when anatomy and urgency permit.
\section*{Risks and Recovery}
Bleeding, infection, injury to nearby organs, blood clots, and lifelong increased risk from encapsulated bacteria are possible. Recommended vaccines and infection precautions are essential.
\section*{References}
\begin{enumerate}\item MedlinePlus. Splenectomy. U.S. National Library of Medicine; 2025.\item Current Medical Diagnosis and Treatment. McGraw-Hill; 2025.\end{enumerate}
\clearpage
